% =====================================================================
% Слайд 17 — Выводы и дальнейшая работа
% =====================================================================
\begin{frame}{Выводы и дальнейшая работа}
  \small

  \textbf{Сделано:}
  \begin{itemize}\footnotesize\setlength{\itemsep}{0.2em}
    \item Систематизировали baseline (OLTW, HMM, HSMM, Hybrid) ---
          все три классических подхода работают на синтетических
          сценариях.
    \item Реструктурировали кодовую базу под расширения, заложили
          типизированные DTO под RL.
    \item Воспроизводимые ноутбуки-демонстрации каждого компонента.
  \end{itemize}

  \textbf{Ключевая идея:}\\
  RL-агент не заменяет классический трекер, а \rlhi{дополняет} его
  как лёгкий слой опережающего предсказания темпа. \\
  \emph{Это новая ниша, в литературе не покрытая.}

  \textbf{Дальше:}
  \begin{itemize}\footnotesize\setlength{\itemsep}{0.2em}
    \item Подключить ASAP, обучить BC $\to$ PPO по описанной схеме.
    \item Сравнить с baseline по RMS render-perf delay и tempo jerk.
    \item Адаптация под конкретного исполнителя (fine-tune политики
          после нескольких сессий) --- главное прикладное
          преимущество RL.
  \end{itemize}

  \textbf{Ограничения:}
  \begin{itemize}\footnotesize\setlength{\itemsep}{0.2em}
    \item Бюджет 20\,мс end-to-end --- нужно подтвердить на крупных
          пьесах с beam-HSMM.
    \item Качество обучения упирается в качество ASAP-разметки и
          симулятора rubato.
    \item Микрофонный вход (требование заказчика) --- отдельная
          подзадача, не покрывается тезисами.
  \end{itemize}
\end{frame}
